% Sperrvermerk ------------------------------------------------------------------
% Nur noetig, wenn die Dokumentation vertrauliche/datenschutzrelevante Daten
% enthaelt (siehe Vorgabe: "Vertrauliche oder datenschutzrelevante Daten sind
% als solche zu kennzeichnen"). Mit Ausbildungsbetrieb und IHK abstimmen. Falls
% nicht benoetigt: Datei leeren oder % Sperrvermerk ------------------------------------------------------------------
% Nur noetig, wenn die Dokumentation vertrauliche/datenschutzrelevante Daten
% enthaelt (siehe Vorgabe: "Vertrauliche oder datenschutzrelevante Daten sind
% als solche zu kennzeichnen"). Mit Ausbildungsbetrieb und IHK abstimmen. Falls
% nicht benoetigt: Datei leeren oder % Sperrvermerk ------------------------------------------------------------------
% Nur noetig, wenn die Dokumentation vertrauliche/datenschutzrelevante Daten
% enthaelt (siehe Vorgabe: "Vertrauliche oder datenschutzrelevante Daten sind
% als solche zu kennzeichnen"). Mit Ausbildungsbetrieb und IHK abstimmen. Falls
% nicht benoetigt: Datei leeren oder % Sperrvermerk ------------------------------------------------------------------
% Nur noetig, wenn die Dokumentation vertrauliche/datenschutzrelevante Daten
% enthaelt (siehe Vorgabe: "Vertrauliche oder datenschutzrelevante Daten sind
% als solche zu kennzeichnen"). Mit Ausbildungsbetrieb und IHK abstimmen. Falls
% nicht benoetigt: Datei leeren oder \input{sperrvermerk} in main.tex entfernen.
\thispagestyle{empty}
\vspace*{4cm}

\begin{center}
  \bfseries Sperrvermerk
\end{center}

\vspace{1cm}

\noindent
[PLATZHALTER: Standardformulierung, z.B. "Die vorliegende
Projektdokumentation enth\"alt vertrauliche Informationen von \dokBetrieb,
insbesondere zu [PLATZHALTER: betroffene Bereiche, z.B. interne Systeme,
Kundendaten, Kennzahlen]. Eine Weitergabe oder Ver\"offentlichung des Inhalts,
auch auszugsweise, ist ohne ausdr\"uckliche schriftliche Zustimmung von
\dokBetrieb nicht gestattet. Ausgenommen hiervon ist die Einsichtnahme durch
die zust\"andigen Pr\"ufungsausschuss- und IHK-Mitglieder im Rahmen der
Abschlusspr\"ufung."]

\vspace{2cm}
\noindent
[PLATZHALTER: Ort, Datum, Unterschrift Ausbildungsbetrieb -- vor Abgabe
erg\"anzen bzw. diese Seite entfernen, falls kein Sperrvermerk gefordert wird.]

\clearpage
 in main.tex entfernen.
\thispagestyle{empty}
\vspace*{4cm}

\begin{center}
  \bfseries Sperrvermerk
\end{center}

\vspace{1cm}

\noindent
[PLATZHALTER: Standardformulierung, z.B. "Die vorliegende
Projektdokumentation enth\"alt vertrauliche Informationen von \dokBetrieb,
insbesondere zu [PLATZHALTER: betroffene Bereiche, z.B. interne Systeme,
Kundendaten, Kennzahlen]. Eine Weitergabe oder Ver\"offentlichung des Inhalts,
auch auszugsweise, ist ohne ausdr\"uckliche schriftliche Zustimmung von
\dokBetrieb nicht gestattet. Ausgenommen hiervon ist die Einsichtnahme durch
die zust\"andigen Pr\"ufungsausschuss- und IHK-Mitglieder im Rahmen der
Abschlusspr\"ufung."]

\vspace{2cm}
\noindent
[PLATZHALTER: Ort, Datum, Unterschrift Ausbildungsbetrieb -- vor Abgabe
erg\"anzen bzw. diese Seite entfernen, falls kein Sperrvermerk gefordert wird.]

\clearpage
 in main.tex entfernen.
\thispagestyle{empty}
\vspace*{4cm}

\begin{center}
  \bfseries Sperrvermerk
\end{center}

\vspace{1cm}

\noindent
[PLATZHALTER: Standardformulierung, z.B. "Die vorliegende
Projektdokumentation enth\"alt vertrauliche Informationen von \dokBetrieb,
insbesondere zu [PLATZHALTER: betroffene Bereiche, z.B. interne Systeme,
Kundendaten, Kennzahlen]. Eine Weitergabe oder Ver\"offentlichung des Inhalts,
auch auszugsweise, ist ohne ausdr\"uckliche schriftliche Zustimmung von
\dokBetrieb nicht gestattet. Ausgenommen hiervon ist die Einsichtnahme durch
die zust\"andigen Pr\"ufungsausschuss- und IHK-Mitglieder im Rahmen der
Abschlusspr\"ufung."]

\vspace{2cm}
\noindent
[PLATZHALTER: Ort, Datum, Unterschrift Ausbildungsbetrieb -- vor Abgabe
erg\"anzen bzw. diese Seite entfernen, falls kein Sperrvermerk gefordert wird.]

\clearpage
 in main.tex entfernen.
\thispagestyle{empty}
\vspace*{4cm}

\begin{center}
  \bfseries Sperrvermerk
\end{center}

\vspace{1cm}

\noindent
[PLATZHALTER: Standardformulierung, z.B. "Die vorliegende
Projektdokumentation enth\"alt vertrauliche Informationen von \dokBetrieb,
insbesondere zu [PLATZHALTER: betroffene Bereiche, z.B. interne Systeme,
Kundendaten, Kennzahlen]. Eine Weitergabe oder Ver\"offentlichung des Inhalts,
auch auszugsweise, ist ohne ausdr\"uckliche schriftliche Zustimmung von
\dokBetrieb nicht gestattet. Ausgenommen hiervon ist die Einsichtnahme durch
die zust\"andigen Pr\"ufungsausschuss- und IHK-Mitglieder im Rahmen der
Abschlusspr\"ufung."]

\vspace{2cm}
\noindent
[PLATZHALTER: Ort, Datum, Unterschrift Ausbildungsbetrieb -- vor Abgabe
erg\"anzen bzw. diese Seite entfernen, falls kein Sperrvermerk gefordert wird.]

\clearpage
